\section{The spectral inequality and its sharp constant}
\label{sec:spectral}

We now apply Theorem~\ref{thm:action} to finite families of negative
eigenfunctions, following the spectral passage of the scalar
companion~\cite[Section~7]{ScalarCompanion}.  The lower triangular coefficients
in that theorem preserve
the spectral bound on each successive span.  A trace version of Young's
inequality then bounds the action by the potential integral.  We first justify
this passage for the measurable potentials in Theorem~\ref{thm:main}.

\subsection{The form domain and negative spectrum}

\begin{lemma}\label{lem:form-compactness}
Let $m\geq1$, let $p>1$, and let $W:\R\to\C^{m\times m}$ be a measurable
Hermitian positive semidefinite function satisfying
$\int_\R\tr(W^p)<\infty$.  The form
\[
 h_W[\psi]=\int_\R\|\psi'\|^2
              -\int_\R\langle\psi,W\psi\rangle,
 \qquad \psi\in H^1(\R;\C^m),
\]
is closed and bounded below.  Multiplication by $W^{1/2}$ is compact from
$H^1(\R;\C^m)$ to $L^2(\R;\C^m)$.  The associated self-adjoint operator has
only isolated eigenvalues of finite multiplicity below zero, with possible
accumulation only at zero.
\end{lemma}

\begin{proof}
Put $w(x)=\|W(x)\|_{\op}$.  Positivity gives
$w^p\leq\tr(W^p)$, so $w\in L^p(\R)$.  The one-dimensional Sobolev estimate
\[
 \|\psi\|_\infty^2\leq2\|\psi\|_2\|\psi'\|_2
\]
holds also for vector-valued functions, by applying the scalar estimate to
their norm.  H\"older's inequality and interpolation therefore give
\begin{align}
 \int_\R\langle\psi,W\psi\rangle
 &\leq\|w\|_p\|\psi\|_{2p/(p-1)}^2 \notag\\
 &\leq 2^{1/p}\|w\|_p
           \|\psi'\|_2^{1/p}\|\psi\|_2^{2-1/p} \notag\\
 &\leq\varepsilon\|\psi'\|_2^2+C_\varepsilon\|\psi\|_2^2
 \qquad(\varepsilon>0).                            \label{eq:form-bound}
\end{align}
The last step is Young's inequality.  Consequently a sufficiently shifted
form norm is equivalent to the $H^1$ norm.  The potential form is continuous
on $H^1$, so the form is closed and bounded below.

For compactness, truncate the multiplier both in space and in size:
\[
 T_{R,L}\psi
   =\mathbf1_{\{|x|\leq R,\ w(x)\leq L\}}W^{1/2}\psi.
\]
Restriction to $[-R,R]$ followed by the compact embedding
$H^1([-R,R];\C^m)\hookrightarrow L^2([-R,R];\C^m)$ and bounded
multiplication shows that $T_{R,L}$ is compact.  If $T\psi=W^{1/2}\psi$,
H\"older's inequality and the same Sobolev embedding imply
\[
 \|(T-T_{R,L})\psi\|_2^2
 \leq C\big\|w\mathbf1_{\{|x|>R\}\cup\{w>L\}}\big\|_p
           \|\psi\|_{H^1}^2.
\]
The coefficient tends to zero as $R,L\to\infty$, proving that $T$ is compact.

Finally, fix $a>0$.  If the spectral subspace of $H_W$ in $(-\infty,-a]$
were infinite-dimensional, it would contain an $L^2$-orthonormal sequence
$(\psi_j)$.  Semiboundedness and the spectral theorem place this sequence in
the form domain with uniformly bounded shifted form norm, hence uniformly
bounded $H^1$ norm.  A weakly convergent subsequence in $H^1$ has limit zero,
because an orthonormal sequence converges weakly to zero in $L^2$.
Compactness of $T$ would then give $\|T\psi_j\|_2\to0$, whereas
\[
 -a\geq h_W[\psi_j]
     =\|\psi_j'\|_2^2-\|T\psi_j\|_2^2
     \geq-\|T\psi_j\|_2^2,
\]
a contradiction.  Every such spectral subspace is therefore finite-dimensional,
which proves the assertion about the negative spectrum.
\end{proof}

\subsection{From the action to spectral moments}

\begin{proof}[Proof of Theorem~\ref{thm:main}]
Recall that $\sigma=\gamma-\tfrac12$, $p=1+\sigma$ and
$q=1+1/\sigma$.  We first prove the inequality for real symmetric potentials.
If there are no negative eigenvalues, it is immediate.  Otherwise fix any
finite nonempty list of negative eigenvalues, respecting their multiplicities.
Since complex conjugation preserves each eigenspace, we may choose corresponding
orthonormal eigenfunctions $\phi_1,\ldots,\phi_n$ that are real-valued.

The action theorem requires functions supported on a bounded interval.  Choose
real smooth compactly supported approximations to the $\phi_i$ in $H^1$.
Their Gram matrices converge to $I_n$, and their form matrices converge to the
diagonal matrix of the selected eigenvalues, by~\eqref{eq:form-bound}.
Multiplying each approximating family by the inverse square root of its Gram
matrix makes its rows orthonormal without changing their common compact
support.  The form matrices still converge to the same diagonal matrix.
For sufficiently good approximations they are negative definite; diagonalize
them by a real orthogonal change of rows.

For one such family, let $U\in H^1_0(\Omega;\R^{n\times m})$ be the matrix
whose rows are the resulting functions $u_i$, where the bounded interval
$\Omega$ contains their supports in its interior.  Extend these rows by zero
when evaluating the whole-line form.  They satisfy
\[
 \int_\Omega UU^{\mathsf T}=I_n,
 \qquad h_W[u_i,u_j]=-k_i^2\delta_{ij},
 \qquad k_1\geq\cdots\geq k_n>0,
\]
where $h_W[\cdot,\cdot]$ is the polarized form and the numbers $-k_i^2$
converge, along the approximations, to the selected eigenvalues in increasing
order.  Set $K=\diag(k_1,\ldots,k_n)$.

Let $C$ be any real lower triangular matrix and write $A=CU$.  Its $i$th row
is $a_i=\sum_{j\leq i}C_{ij}u_j$, so
\begin{equation}\label{eq:triangular-spectral-order}
 h_W[a_i]+k_i^2\|a_i\|_{L^2}^2
   =\sum_{j\leq i}C_{ij}^2(k_i^2-k_j^2)\leq0.
\end{equation}
Summing over $i$ bounds the quadratic terms of the action.  Since
$AA^{\mathsf T}$ and $A^{\mathsf T}A$ have the same nonzero eigenvalues,
\begin{equation}\label{eq:spectral-action-upper}
 \mathcal E(CU)
 \leq\int_\Omega
      \bigl\{\tr(WA^{\mathsf T}A)-\tr((A^{\mathsf T}A)^q)\bigr\}\,dx.
\end{equation}

The right side has a pointwise bound that does not require its two matrices to
commute.  Indeed, scalar maximization gives, for $y,z\geq0$,
\begin{equation}\label{eq:scalar-young-sharp}
 yz-z^q\leq D_\sigma y^p,
 \qquad D_\sigma=\frac{\sigma^\sigma}{(1+\sigma)^{1+\sigma}}.
\end{equation}
To apply it, diagonalize an arbitrary positive semidefinite matrix $Z$ in an
orthonormal basis $(v_j)$, with eigenvalues $z_j$.  Set
$y_j=\langle v_j,Wv_j\rangle$.  Convexity of $t\mapsto t^p$, applied to the
spectral measure of $W$ in $v_j$, gives
$y_j^p\leq\langle v_j,W^pv_j\rangle$.  Hence
\begin{equation}\label{eq:trace-young-sharp}
 \tr(WZ)-\tr(Z^q)
   =\sum_j(y_jz_j-z_j^q)
   \leq D_\sigma\sum_j y_j^p
   \leq D_\sigma\tr(W^p).
\end{equation}
Using $Z=A^{\mathsf T}A$ in~\eqref{eq:spectral-action-upper}, we obtain
\[
 \sup_{C\ \mathrm{lower\ triangular}}\mathcal E(CU)
    \leq D_\sigma\int_\R\tr(W^p)\,dx.
\]
All terms are finite for each $C$: the rows of $A$ are bounded and supported
on $\overline\Omega$, and $W$ is locally integrable.

Theorem~\ref{thm:action} supplies the reverse bound in terms of the $k_i$.
Writing $\mathrm B$ for Euler's beta function, the substitution $s=k_it$ gives
\[
 \int_{-k_i}^{k_i}(k_i^2-s^2)^\sigma\,ds
  =k_i^{2\sigma+1}\mathrm B(\tfrac12,1+\sigma)
  =k_i^{2\gamma}\mathrm B(\tfrac12,1+\sigma).
\]
We conclude that
\begin{equation}\label{eq:finite-spectral-moment}
 \sum_{i=1}^n k_i^{2\gamma}
   \leq\frac{D_\sigma}{\mathrm B(\tfrac12,1+\sigma)}
            \int_\R\tr(W^p)\,dx.
\end{equation}
The coefficient is exactly the one in the theorem:
\begin{align}
 \frac{D_\sigma}{\mathrm B(\tfrac12,1+\sigma)}
 &=\frac{\sigma^\sigma}{(1+\sigma)^{1+\sigma}}
       \frac{\Gamma(\sigma+3/2)}{\sqrt\pi\,\Gamma(\sigma+1)} \notag\\
 &=\left(\frac{\sigma}{1+\sigma}\right)^\sigma
       \frac{\Gamma(\sigma+3/2)}{\sqrt\pi\,\Gamma(\sigma+2)}
   =C_\gamma.                                           \label{eq:constant-identity}
\end{align}
Letting the approximations converge proves the bound for the chosen finite
list of eigenvalues.  The sum of all negative eigenvalue moments is the
supremum of these finite sums.  Thus the same bound holds for
$\Tr(H_W)_-^\gamma$, and in particular this trace is finite.

For a complex Hermitian potential, write $W=E+iF$ with real matrices $E,F$.
Hermiticity means $E^{\mathsf T}=E$ and $F^{\mathsf T}=-F$, so
\[
 \widetilde W=\begin{pmatrix}E&-F\\ F&E\end{pmatrix}
\]
is real symmetric.  The constant complex unitary
\[
 \mathcal U=\frac1{\sqrt2}
             \begin{pmatrix}I&iI\\ I&-iI\end{pmatrix}
\]
satisfies $\mathcal U\widetilde W\mathcal U^*
=W\oplus\overline W$.  It follows that $\widetilde W\geq0$ and
$\tr(\widetilde W^p)=2\tr(W^p)$.  Since $\mathcal U$ also commutes with
differentiation, the same change of coordinates gives
\[
 H_{\widetilde W}\simeq H_W\oplus H_{\overline W}.
\]
Complex conjugation intertwines the two summands and preserves eigenvalue
multiplicities.  The spectral moment for $\widetilde W$ is consequently twice
that for $W$.  Applying the real inequality in dimension $2m$ and dividing
both sides by two proves the asserted inequality for all Hermitian $W$.

It remains to verify sharpness.  Put $r=1/\sigma>1$ and consider the scalar
potential and function
\[
 V(x)=(r+1)\sech^2(rx),
 \qquad \psi(x)=\sech^{1/r}(rx).
\]
The identity $\psi'/\psi=-\tanh(rx)$ gives $H_V\psi=-\psi$, and
$\psi\in L^2(\R)$.  We determine the full negative spectrum by the standard
factorization method; see~\cite[Sections~II--III]{BenguriaLoss2000} for its use
in Lieb--Thirring inequalities.  Define
$Q=\partial_x+\tanh(rx)$ on $H^1(\R)$.  Direct multiplication yields
\[
 H_V=Q^*Q-1,
 \qquad QQ^*-1=-\partial_x^2+(r-1)\sech^2(rx)\geq0.
\]
Thus $H_V\geq-1$, and its eigenspace at $-1$ is $\ker Q$, which is spanned
by $\psi$.  If $\lambda\in(-1,0)$ were an eigenvalue with eigenfunction $u$,
then $Qu\ne0$.  The bounded potential gives $u\in H^2$, so $Qu\in L^2$.
The intertwining identity $(QQ^*-1)Qu=\lambda Qu$ holds in distributions.
Boundedness of the partner potential then implies $Qu\in H^2$, so this
identity contradicts nonnegativity.  Therefore $-1$ is the only negative
eigenvalue and is simple.
Finally,
\[
 \int_\R V^p\,dx
  =\frac{(r+1)^p}{r}\,\mathrm B(\tfrac12,p)
  =\frac{\mathrm B(\tfrac12,1+\sigma)}{D_\sigma}
  =C_\gamma^{-1}.
\]
The scalar inequality is an equality for $V$.  Embedding $V$ in one diagonal
channel and setting the other channels to zero gives equality in every
matrix dimension, proving that $C_\gamma$ is optimal.
\end{proof}
